\chapter{Results and Discussion}
\label{ch:results}

% DRAFT, 2026-09-16. Sections 4.1-4.4 report measured results and are complete for the reference
% controllers; Section 4.5 is a placeholder until the MARL training run and its evaluation finish.
% Every number is reproducible from starter/results/ (the source file is named beside each table).
% This chapter stays commented out of main.tex until Section 4.5 is filled in.

This chapter reports what the constructed environment reproduces, how the three reference controllers
behave in it, and how that behaviour degrades as disturbances are activated. Unless stated otherwise,
every cell is $N = 30$ paired Monte Carlo replications: a seed fixes one realization of demand, dwell,
running time, dispatch deviation, surge, weather, and breakdown, and every controller faces that same
realization, so differences between controllers are attributable to control alone. Intervals are
bootstrap 95\% confidence intervals from 5{,}000 resamples.

\section{Environment Construction and Validation (EO 1.1)}
\label{sec:results-eo11}

The corridor is the 27 modelled stops of Route 801 direction code 6, built on the real road geometry
taken from the OpenStreetMap route relation. Segment running times, dwell parameters, per-trip demand,
and dispatch deviation are estimated from the 229{,}421 clean APC stop events, using calibration
service days only.

\subsection{Segment running time}

Edge speeds were tuned on the calibration days and then evaluated once on the held-out days.

\begin{table}[htbp]
\centering
\caption[Corridor calibration]{Corridor calibration against APC running times, 26 segments.}
\label{tab:calibration}
\begin{tabular}{lccc}
\toprule
\textbf{Service days} & \textbf{RMSPE} & \textbf{$GEH < 5$} & \textbf{Max GEH} \\
\midrule
Calibration (65 days, alternating)     & 0.90\% & 26/26 & 0.25 \\
Held out (64 days, alternating)        & 3.08\% & 26/26 & 0.90 \\
Held out (chronological split, check)  & 6.83\% & 26/26 & 3.59 \\
\bottomrule
\end{tabular}
\end{table}

The environment satisfies the acceptance criterion of Section~\ref{sec:research-design} on days it
never saw. The third row is a robustness check rather than the study's split: when service days are
divided chronologically instead of alternately, the held-out error doubles while every segment still
satisfies $GEH < 5$. Recorded weekday boardings rise about 31\% from July to October while the median
segment running time moves by less than 5\%, so that figure reflects a seasonal shift in the held-out
period rather than a worse model.

\subsection{Headway variability, load, and stop service}

Travel-time calibration alone does not establish that the environment reproduces \emph{bunching}, which
is the quantity this study reports on. Three further checks compare the uncontrolled environment with
the held-out days, each measured identically on both sides.

\begin{table}[htbp]
\centering
\caption[Environment validation]{Uncontrolled environment against held-out observations.}
\label{tab:validation}
\begin{tabular}{lcc}
\toprule
\textbf{Quantity} & \textbf{Observed} & \textbf{Simulated (No Control)} \\
\midrule
Headway CV, corridor mean              & 0.504 & 0.556 \\
Headway CV, first stop                 & 0.352 & 0.357 \\
Headway CV, stop-by-stop correlation   & \multicolumn{2}{c}{$r = 0.87$} \\
Share of trips serving a stop          & 0.64  & 0.69 \\
Stop-service correlation               & \multicolumn{2}{c}{$r = 0.96$} \\
Riders on board, profile correlation   & \multicolumn{2}{c}{$r = 0.99$} \\
\bottomrule
\end{tabular}
\end{table}

Dispatch variability is anchored to the observed first-stop headway CV, which the environment
reproduces to within 0.005. Bunching then accumulates about 10\% faster than observed over the second
half of the corridor. Real operators hold buses at timepoints, which would damp exactly that
accumulation; the APC schedule field is populated on only 18\% of records, so this explanation cannot
be tested with the available data and the residual is reported as a limitation rather than tuned away.
Simulated loads sit about two riders below the mean APC \texttt{max\_load} because APC records a stop
only when its doors open, which over-samples busier trips; the shape along the corridor matches.

\section{Disturbance Generators (EO 1.2)}
\label{sec:results-eo12}

The environment exposes each disturbance class as an adjustable parameter, and every value below is
either estimated from the data or declared synthetic.

\begin{table}[htbp]
\centering
\caption[Disturbance settings]{Disturbance generators and their evaluation settings.}
\label{tab:disturbance-settings}
\begin{tabular}{llp{7.2cm}}
\toprule
\textbf{Class} & \textbf{Parameter} & \textbf{Setting} \\
\midrule
D & dispersion 3.8   & Negative-binomial demand; variance-to-mean ratio estimated from the APC counts \\
T & fitted spreads   & Running-time and dispatch spreads estimated from neighbouring-bus pairs \\
S & $\sigma_d = 1$   & Demand multiplier $f_d \sim N(1, \sigma_d^2)$ clipped to $[1, 10]$ \\
W & rain $\times \eta$ & Observed ordinary-rain multiplier 1.0135 (not significant), times a labelled synthetic lognormal stress of coefficient of variation $\eta \in \{0, 0.3, 0.6, 1.0, 1.3\}$ \\
B & $n_B = 1$        & One bus removed from service at a random stop; its riders alight and wait for the next bus \\
\bottomrule
\end{tabular}
\end{table}

The ordinary-rain effect is small and statistically indistinguishable from zero: running times on
rain-exposed records are 1.35\% longer within segment, time-of-day, and day-type strata, with a 95\%
interval of $[-0.8\%, +2.5\%]$ across 170 strata. That is reported as the observed weather condition,
and anything stronger is labelled synthetic throughout.

\section{Reference Controller Performance (EO 3.1)}
\label{sec:results-eo31}

\begin{table}[htbp]
\centering
\caption[Stage A results]{Stage A (D+T), $N = 30$ paired seeds. Source: \texttt{results/mc\_summary.md}.}
\label{tab:stage-a}
\begin{tabular}{lccc}
\toprule
\textbf{Controller} & \textbf{Headway CV} & \textbf{Wait (s)} & \textbf{Travel (s)} \\
\midrule
No Control      & 0.556 [0.518, 0.590] & 387 [375, 400] & 4391 \\
Forward Headway & 0.396 [0.372, 0.420] & 350 [342, 358] & 4533 \\
Even Headway    & \textbf{0.342} [0.322, 0.363] & \textbf{339} [333, 346] & 4543 \\
\bottomrule
\end{tabular}
\end{table}

Both holding rules reduce bunching significantly on an ordinary day: Forward Headway by 29\%
[22, 35] and Even Headway by 38\% [33, 44]. Mean passenger waiting falls by 10\% and 12\%
respectively, while in-vehicle travel time rises by about 3\%, which is the holding time itself. Even
Headway is the stronger reference throughout, as the literature reports: it observes both neighbouring
gaps, whereas Forward Headway is the deliberately gentle rule of Daganzo with gain $\alpha = 0.2$.

\section{Performance Under Non-Ideal Conditions (EO 3.2)}
\label{sec:results-eo32}

\begin{table}[htbp]
\centering
\caption[Single-disturbance ablations]{Headway CV by activation-matrix condition, $N = 30$ paired
seeds. The change is measured against No Control in the same condition.}
\label{tab:ablations}
\begin{tabular}{lccc}
\toprule
\textbf{Condition} & \textbf{NC} & \textbf{FH} & \textbf{EH} \\
\midrule
Stage A (D+T)                     & 0.556 & 0.396 ($-29\%$) & 0.342 ($-38\%$) \\
S ablation (D+T+S)                & 0.596 & 0.427 ($-28\%$) & 0.371 ($-38\%$) \\
W ablation (D+T+W, observed rain) & 0.559 & 0.400 ($-28\%$) & 0.346 ($-38\%$) \\
B ablation (D+T+B)                & 0.564 & 0.423 ($-25\%$) & 0.370 ($-34\%$) \\
Stage B (all, observed rain)      & 0.609 & 0.458 ($-25\%$) & 0.404 ($-34\%$) \\
\bottomrule
\end{tabular}
\end{table}

Under observed conditions the reference controllers are robust. A demand surge raises bunching for
every controller but leaves the holding advantage unchanged; a breakdown costs each rule about four
percentage points of its advantage; and observed ordinary rain is nearly indistinguishable from
Stage A, as the estimated rain effect implies. With all three active, Even Headway still cuts bunching
by about a third.

\begin{table}[htbp]
\centering
\caption[Stage B weather sweep]{Stage B (D+T+S+W+B) as the labelled synthetic weather stress $\eta$
increases. Source: \texttt{results/stageB\_weather\_sweep.csv}.}
\label{tab:stage-b-sweep}
\begin{tabular}{lcccc}
\toprule
\textbf{Weather} & \textbf{NC} & \textbf{FH} & \textbf{EH} & \textbf{EH wait (s)} \\
\midrule
Observed rain only & 0.609 & 0.458 ($-25\%$) & 0.404 ($-34\%$) & 361 \\
$\eta = 0.3$       & 0.685 & 0.574 ($-16\%$) & 0.534 ($-22\%$) & 399 \\
$\eta = 0.6$       & 0.820 & 0.742 ($-9\%$)  & 0.725 ($-12\%$) & 472 \\
$\eta = 1.0$       & 0.904 & 0.836 ($-8\%$)  & 0.832 ($-8\%$)  & 524 \\
$\eta = 1.3$       & 0.929 & 0.862 ($-7\%$)  & 0.864 ($-7\%$)  & 541 \\
\bottomrule
\end{tabular}
\end{table}

Degradation is driven almost entirely by weather stress. As $\eta$ increases, the advantage of holding
shrinks monotonically, and beyond $\eta = 1.0$ the two rules are statistically indistinguishable from
each other: local holding cannot restore spacing once the disturbance to running time is of the same
order as the headway itself. That is the regime in which an adaptive controller has something to gain,
and it lies outside the observed weather support, so it is reported as a labelled synthetic
extrapolation and never as a rainfall category.

\section{MARL Controller Performance}
\label{sec:results-marl}

% PENDING: the domain-randomized training run (D+T always active; S, W and B randomized per episode)
% and the evaluation of its frozen policy over the nine-cell matrix. The acceptance criteria and the
% statistical procedure were fixed before that run and are stated in Section 3.5 and in
% scripts/eval_marl.py:
%   Stage A (i)  mean wait not significantly worse than Even Headway
%   Stage A (ii) headway CV significantly below No Control
%   Stage B      mean wait below the best baseline in every Stage B cell
%   plus the project's own training gate: Stage A headway CV below Even Headway (0.342)
% Tests: paired Wilcoxon signed-rank by seed, Holm-corrected across the three baseline comparisons,
% alpha = 0.05, seeds 0-29 for every controller.
